当我们只知道“已有一个波前”，为什么仍能预言下一时刻的新波前会在哪里？若把整个波前当成一块刚性板，便很难解释波为何会绕过障碍物、在边缘重新铺开。换一个角度看，波前上每一点都像一个微小的新波源，都会继续把振动向外传。下一时刻的新波前位置不再由旧波前整体的“平移”决定，而由这些子波共同包出的新外沿确定。

\begin{theorem}[惠更斯原理]
波阵面上的每一点都可以看成发出子波的新波源；经过一段时间后，这些子波的包络面就是新的波前。
\end{theorem}

当波前来到孔缝或障碍物边缘时，继续向前的是边缘附近仍能传播的那一段波前，并不是一条刚性的几何边界。按惠更斯观点，这段波前上的各点仍会发出子波。阴影区并不是“完全没有后续波源”：它会收到这些边缘子波的叠加。

\begin{theorem}[波的反射和折射]
反射与折射的基本规律可并排记成
\[
\begin{aligned}
\text{反射：}\quad &\text{入射线、反射线和界面法线共面},\qquad \text{反射角}=\text{入射角},\\
\text{折射：}\quad &\text{入射线、折射线和界面法线共面},\qquad \frac{\sin i}{\sin r}=\frac{u_1}{u_2}.
\end{aligned}
\]
这条折射关系很值得直观理解：如果第二种介质中波速更小，即 \(u_2<u_1\)，那么折射线会向法线偏折；若 \(u_2>u_1\)，则会偏离法线。
\end{theorem}
